% !TEX program = pdflatex
% ECT UDG stress test: consolidated method --- v4 (post-GPT-round2).
%
% Version 4 changes vs v3 (applying GPT's second-round critique):
%   - Removed "Data cannot falsify theory, only closure" philosophical shortcut
%   - S3 FCC 224: demoted numerical factor to order-unity proxy claim only
%   - S1 UFDs: demoted from "validation" to "side diagnostic"
%   - S2 GCs: kept as Newton sanity check (no change in status)
%   - S4 G_eff(X): substantial weakening. No longer "substantive finding".
%     Now a "candidate theory-level direction" only; no derivation connects
%     preprint's G_eff(phi) to practical galactic closure for DF4/FCC 224.
%   - Removed "not new phenomenology" rhetorical claim (unsupported)
%   - Removed dimensional plausibility hand-wave for delta-phi ~ -ln(10)/beta
%   - Honest classification (III): G_eff path listed equal-footing with others
%   - Conclusion: G_eff not singled out as "most promising"; all candidate
%     routes listed neutrally
%   - Abstract: removed claim that G_eff "is structurally present in preprint"
%
% This document is INTERNAL WORKING NOTES, not a subsection for publication.

\documentclass[11pt]{article}
\usepackage[margin=2.2cm]{geometry}
\usepackage{amsmath,amssymb,amsthm}
\usepackage{mathtools}
\usepackage{hyperref}
\usepackage{booktabs}
\usepackage{graphicx}
\usepackage{longtable}
\hypersetup{colorlinks=true,linkcolor=blue,citecolor=blue,urlcolor=blue}

\title{DM-deficient UDGs as a stress test of the current ECT Level-B\\
galactic closure --- consolidated diagnostic notes}
\author{ECT working notes --- v4 consolidated (internal)}
\date{2026-04-20}

\begin{document}
\maketitle

\begin{abstract}
This v4 supersedes v1 (which incorrectly promoted
$\bar\phi_{\rm local}\to 0^-\Rightarrow g^\dagger_{\rm eff}\to 0$ to
a derived consequence; that inverted the preprint's own regime
ordering at \S44487--44508) and v2--v3 (honest in structure but still
overstated the status of the intra-preprint $G_{\rm eff}(X)$ route).
The present version (i) integrates a three-level framework D/C/T
separating Data, Closure test, and Theory interpretation; (ii) supplies
the inverse-problem closed form
$\Xi_{\rm req}=r(r-1)\,g_N/g^\dagger_0$ with its correct evaluation at
$R_{1/2}$; (iii) documents nuisance and uncertainty budget in which the
stress-test signature is robust; (iv) notes that the preprint's full
field equations contain a dynamical $G_{\rm eff}(\phi)$ which makes a
suppression route structurally \emph{conceivable} but --- critically
--- \emph{no derivation in the attached preprint maps that degree of
freedom onto the present galactic closure for DF4/FCC 224}; (v) lists
secondary factors as future work; and (vi) provides a minimal
integration proposal for the preprint. The scientific bottom line:
DF4 and probably FCC 224 remain sharp stress tests of the current
Level-B galactic closure, not explained consequences of it.
\end{abstract}

\section{Framework: three levels D, C, T}

Three levels must be kept strictly separate.

\paragraph{Level D (Data).} What is directly measured for each object:
$\sigma_{\rm obs}$ (or $v_{\rm rot}$), $M_\star$, $R_e$, $D$, $n$
(S\'ersic), environment proxies (host mass, projected distance), tracer
family (stars vs GCs vs HI), aperture $R_{\rm ap}$, quality flag
$Q_{\rm eq}\in\{{\rm relaxed, possibly\_disturbed, disturbed}\}$.

\paragraph{Level C (Closure test).} What the current practical galactic
closure predicts: for each object, compute the required $\Xi_{\rm req}$
such that the closure reproduces observed kinematics. Compare to the
expected range $\Xi\sim O(1)$ for normal galaxies on the RAR.

\paragraph{Level T (Theory interpretation).} What Level-C results imply
for ECT: if $\Xi_{\rm req}$ extreme, the closure is strained at the
object; this does \emph{not} falsify the full ECT field equations
[eq.~17251]; only the current Level-B closure is directly tested.

\paragraph{Honesty boundary.} In the present context the data directly
test the current Level-B galactic closure; they constrain the broader
theory only through whatever future derivation is able to connect the
full field equations to that closure. No such derivation is given in
the attached preprint.

\section{Base equations (preprint anchors)}

\begin{enumerate}
\item[{[A1]}] Full time-dependent IR action [\S17224]:
$$
S_{\rm IR}=\int d^4x\sqrt{-g}\left[\tfrac{\bar M_{\rm Pl}^2}{2}e^{\beta\phi}R
-\tfrac{\bar M_{\rm Pl}^2}{2}\omega(\phi)(\partial_\mu\phi)^2 - U(\phi)+\mathcal L_{\rm mat}\right].
$$

\item[{[A2]}] Generalized Einstein eqs [\S17248] with dynamical
$G_{\rm eff}(\phi)$, $\Lambda_{\rm eff}(\phi)$ and orientation term
$\alpha_{\rm ECT}(n_\mu n_\lambda R^\lambda_\nu + \ldots)$.

\item[{[A3]}] Quasi-static local $\phi$-equation [\S22486]:
$K(\phi)\Box_{\rm static}\phi \approx (8\pi G_N/c^4)\rho_{\rm bar}\xi_1(\phi)$.

\item[{[A4]}] Practical galactic closure [\S23519]:
$$
g(R)=\tfrac{1}{2}\left[g_N(R)+\sqrt{g_N^2(R)+4g_N(R)\,g^\dagger_{\rm eff}}\right]
$$
with limits $g\to g_N$ (Newton, large $g_N$) and $g\to\sqrt{g_N g^\dagger_{\rm eff}}$
(deep-MOND, small $g_N$). Exactly two asymptotic branches; no third
independent low-$g$ Newtonian branch exists within this closure.

\item[{[A5]}] Environment-modulated critical scale [\S22873]:
$g^\dagger_{\rm eff} = g^\dagger_0\,\Xi(\rho_{\rm env}, \bar\phi_{\rm env},\ldots)$,
with baseline $g^\dagger_0 = cH_0/(2\pi)\approx 1.04\times 10^{-10}\,\text{m/s}^2$.

\item[{[A6]}] Level-B minimal ansatz [\S23718]:
$\Xi = 1/\sqrt{1+\rho_{\rm env}/\rho_*}$. Preprint explicitly flags this
as Level-B interpolation, not a first-principles derivation; exact
$\Xi$-law remains open (OP3) [\S51895].

\item[{[A7]}] Regime ordering [\S44487]: in the preprint's own notation
the nonlinear branch activates at small local $g_N$, not via
$\bar\phi_{\rm local}\to 0^-$ as a separate trigger for
$g^\dagger_{\rm eff}\to 0$. The v1 draft inverted this; v4 respects it.

\item[{[A8]}] SPARC-fitted distribution of $\Xi$: median 0.86, MW 1.72.
DDO 154 (void dwarf) reproduced at baseline [\S23325].
\end{enumerate}

\section{Method: closed form and forward simulation}

\subsection{Closed-form inverse: $\Xi_{\rm req}$ from $\sigma_{\rm obs}$}

Inverting [A4] at $R=R_{1/2}=(4/3)R_e$ with Wolf-type mass estimator:
$$
r \equiv \frac{M_{\rm dyn}}{M_{\rm bar}}\bigg|_{R_{1/2}}
= \frac{g(R_{1/2})}{g_N(R_{1/2})}
= \tfrac{1}{2}\left[1+\sqrt{1+4y}\right], \quad y\equiv\frac{g^\dagger_{\rm eff}}{g_N}
$$
$$
\boxed{\;y = r(r-1), \qquad \Xi_{\rm req} = r(r-1)\cdot\frac{g_N(R_{1/2})}{g^\dagger_0}\;}
$$

Remaining phenomenology:
\begin{itemize}
\item[{[PHEN-1]}] Wolf estimator $M_{\rm dyn,1/2} = 4\sigma^2 R_{1/2}/G$:
isotropic-Plummer proxy, uncertainty factor $\sim 1.3$.
\item[{[PHEN-2]}] $M_{\rm bar,1/2}=f_e M_\star$ with $f_e\approx 0.5$:
$\pm 20\%$ profile uncertainty.
\end{itemize}
Further reduction requires Jeans with explicit profile (\S~\ref{sec:jeans}).

\subsection{Jeans forward model (reduction, not elimination, of phenomenology)}
\label{sec:jeans}

Isotropic spherical Jeans equation:
$$
\sigma_r^2(r) = \frac{1}{\nu(r)J_\beta(r)}\int_r^\infty \nu(s)\,g(s)\,J_\beta(s)\,ds,
\quad
J_\beta(r) = \exp\!\left(\int^r\!\!\frac{2\beta(u)}{u}\,du\right),
$$
with $\nu(r)$ the deprojected tracer density and $g(r)$ from [A4] at
trial $\Xi$. LOS projection and aperture averaging as standard:
$$
\sigma_{\rm los}^2(R)=\frac{2}{I(R)}\int_R^\infty\!
\left(1-\beta\tfrac{R^2}{r^2}\right)\frac{\nu(r)\sigma_r^2(r)\,r\,dr}{\sqrt{r^2-R^2}},
$$
$$
\langle\sigma_{\rm los}^2\rangle_{\rm ap}
= \frac{\int_0^{R_{\rm ap}}W(R)\,I(R)\,\sigma_{\rm los}^2(R)\,2\pi R\,dR}
{\int_0^{R_{\rm ap}}W(R)\,I(R)\,2\pi R\,dR}.
$$

For DF4 (isotropic Plummer, $M/L=2$, $D=20$~Mpc): direct Jeans returns
$\sigma_N^{\rm Jeans}=7.97$~km\,s$^{-1}$ vs Wolf-proxy 6.15~km\,s$^{-1}$
(ratio 1.30). The qualitative stress-test conclusion is unchanged:
$\Xi_{\rm req}\ll 1$ required for DF4 regardless of proxy vs Jeans.

\subsection{Benchmark comparison grid}

For each object, compute predictions under four independent benchmarks
and never mix them:
$$
\begin{array}{ll}
\text{N:}&\text{Newton, } g=g_N\\
\text{M0:}&\text{MOND-AQUAL at }a_0=1.2\times 10^{-10}\,\text{m s}^{-2}\\
\text{ME:}&\text{MOND with EFE modification}\\
\text{E0:}&\text{Current ECT closure [A4] at Level-B }\Xi\text{ [A6]}
\end{array}
$$
Of these, only E0 is ECT; N/M0/ME are astrophysics-community references.

\subsection{Posterior on $\Xi$ (correct final output)}

Per object:
$$
P(\Xi\mid\text{data}) \propto \int d\theta\,\mathcal{L}(\text{data}\mid\Xi,\theta)\,\pi(\Xi)\,\pi(\theta)
$$
with $\theta = \{M/L, D, \beta(r), n, R_{\rm ap}, f_{\rm int}, M_{\rm host}, d_{3D}, Q_{\rm eq}\}$.
Distance $D$ sampled jointly with $M_\star\propto D^2$ and $R_e\propto D$
covariantly. Small-$N$ tracers: individual-velocity likelihood
$\mathcal{L}=\prod_i\mathcal{N}(v_i\mid\bar v,\sqrt{\sigma_m^2+\delta_i^2})$
with optional interloper mixture.

\section{Numerical diagnostic: six objects}

Using Wolf-based $\Xi_{\rm req}$ at $R_{1/2}$, central values with
nuisance bands summarised separately (\S~\ref{sec:nuisance}):

\begin{longtable}{@{}l r r r r r l@{}}
\toprule
Object & $M_\star/10^8M_\odot$ & $R_e$/kpc & $\sigma_{\rm obs}$/km\,s$^{-1}$ &
$M_{\rm dyn}/M_{\rm bar}$ & $\Xi_{\rm req}$ & Class\\
\midrule
\endhead
\bottomrule
\endfoot
DF4           & 1.5  & 1.6  & $6.3^{+2.5}_{-1.6}$ & $1.1$   & $\sim 10^{-3}$      & III (stress-low) \\
FCC 224       & 1.74 & 1.89 & $7.8^{+6.7}_{-4.4}$ & $1.6$   & $\sim 2\times 10^{-2}$ & III (stress-low) \\
DF2           & 1.3  & 2.2  & $[8.6,\,14.9]$      & $3$--$9$ & $0.1$--$1.4$        & borderline \\
NGC 5846-UDG1 & 1.1  & 2.1  & $17\pm 2$           & $13.7$  & $\sim 1.6$          & II (normal RAR) \\
Dragonfly 44  & 3.0  & 4.7  & $33\pm 3$           & $42.3$  & $\sim 8.9$          & II (upper tail) \\
AGC 114905    & \multicolumn{5}{l}{Gas-rich rotator; spherical Jeans not applicable; disc solver required}& II$'$ \\
\end{longtable}

Matched-pair observation: DF4 and NGC~5846-UDG1 have comparable
$M_\star\sim 10^8M_\odot$, $R_e\sim 2$~kpc, and comparable group-scale
environment, yet their inferred $\Xi_{\rm req}$ differ by $\sim 10^3$.

\subsection{Side diagnostics and sanity checks}

\paragraph{S1. Ultra-faint dwarfs (side diagnostic, not validation).}
For $M_\star\sim 10^3$--$10^7\,M_\odot$ (Draco, Sculptor, Fornax,
Segue~I, Willman~I), the Wolf-based method returns $\Xi_{\rm req}$
anywhere from $\sim 2$ to $\sim 10^3$. This is not a ``method works''
result --- it is simply the statement that the Wolf-proxy chain does not
obviously collapse outside the UDG target mass range. The spread is
driven by known, well-documented effects (binary-star $\sigma$
inflation, tidal non-equilibrium, radial anisotropy, small-$N$
statistics) that are orthogonal to the UDG stress-test conclusion.
Listed here only for context.

\paragraph{S2. Globular clusters (Newton sanity check).}
47~Tuc, $\omega$~Cen, M13, M15: $g_N/g^\dagger_0 \gtrsim 15$, placing
them safely in the $g_N\gg g^\dagger_{\rm eff}$ branch of [A4]. The
closure returns $g\to g_N$ automatically, independent of $\Xi$. This is
the trivial Newton-recovery check and is listed for completeness.

\paragraph{S3. FCC 224 non-sphericity.}
FCC~224 is prolate with $\varepsilon\simeq 0.36$. Spherical Jeans
should therefore be treated only as an order-unity proxy for this
object; no reliable numerical correction factor can be claimed from
the attached preprint alone. The central $\Xi_{\rm req}$ value for
FCC~224 should be read as order-of-magnitude, not precision. This does
not change the qualitative classification of FCC~224 as stress-low.

\paragraph{S4. $G_{\rm eff}(\phi)$: a candidate theory-level direction only.}
The generalized field equations [A2] contain a dynamical $G_{\rm eff}(\phi)$,
which in principle makes a suppression route structurally conceivable.
In a pure forward parameter scan at fixed $\Xi=1$, simply replacing
$G_N$ by $\eta\,G_N$ with $\eta\in[0.1,0.3]$ reduces the predicted
$\sigma$ for DF4 from $\sim 19$~km\,s$^{-1}$ to $\sim 6$--$10$~km\,s$^{-1}$,
which would fall in the observed band.
\textbf{This is not a derivation.} The preprint does not contain a
derivation that maps the $G_{\rm eff}(\phi)$ degree of freedom of
eqs.~[A1]--[A2] onto the current practical galactic closure [A4] for
DF4/FCC~224: [A4] is written in terms of $g_N(R)$ and
$g^\dagger_{\rm eff}$, not in terms of $G_{\rm eff}(\phi)$. Promoting
the parameter scan above to an ``intra-ECT rescue mechanism'' would
therefore require an additional connecting derivation, which is not
available in the attached files.

\paragraph{S5. Selection rate.}
Approximately 3 confirmed DM-deficient UDGs out of $\sim 500$ screened
($\sim 0.6\%$), consistent in order of magnitude with the $\sim 1\%$
rate predicted by $\Lambda$CDM bullet-dwarf scenarios [Moreno+2022].
Any future ECT-compatible rescue mechanism will need to be consistent
with this observed rate.

\section{Nuisance and uncertainty budget}
\label{sec:nuisance}

For DF4, total $\Xi_{\rm req}$ uncertainty across independent nuisances:

\begin{center}
\begin{tabular}{ll}
\toprule
Source & Multiplicative factor on $\Xi_{\rm req}$\\
\midrule
Distance $D\in[13,25]$ Mpc (covariant) & $\sim 10\times$\\
$M/L\in[1,4]$ & $\sim 10\times$\\
Small-$N$ at 95\% CL on $\sigma$ & $\sim 100\times$\\
Anisotropy $\beta\in[-0.5,0.5]$ & $\sim 1.5\times$\\
Proxy $\to$ Jeans & $\sim 1.3\times$\\
\bottomrule
\end{tabular}
\end{center}

Combined 95\% CL: $\log_{10}\Xi_{\rm req}\in[-4,\,-0.5]$ for DF4. The
entire interval lies below the RAR-normal range $\Xi\sim O(1)$. The
stress-test classification is therefore robust to nuisance variations;
only the numerical value of $\Xi_{\rm req}$ is imprecise.

\section{Structural inadequacy of the density-only $\Xi$-ansatz}

This is a structural observation, not an impossibility theorem.

The Level-B ansatz [A6] $\Xi = 1/\sqrt{1+\rho_{\rm env}/\rho_*}$ gives
$\Xi\leq 1$ for any $\rho_{\rm env}$, and depends monotonically on a
single environmental scalar. For the matched pair DF4 vs
NGC~5846-UDG1:
\begin{itemize}
\item comparable $M_\star,R_e$ (both $\sim 10^8 M_\odot$, $\sim 2$ kpc),
\item comparable group environment,
\item inferred $\Xi_{\rm req}$ differing by $\sim 10^3$.
\end{itemize}
A density-only ansatz would need a $\rho_{\rm env}$ contrast of order
$10^6$ to accommodate this --- not observed. Therefore the minimal
Level-B ansatz [A6] is \emph{structurally insufficient} for this class
of objects. This does not rule out ECT-compatible environment laws in
general; the preprint explicitly allows more general functional forms
[\S23722] as part of the open problem OP3.

\section{EFE benchmark (external to ECT)}

As a benchmark parallel to the ECT analysis, the MOND-with-EFE
expectation gives, for the four pressure-supported systems:

\begin{center}
\begin{tabular}{lccc}
\toprule
Object & $M_{\rm host}/M_\odot$ & $d_{\rm proj}$/kpc & $g_{\rm ext}/a_0$\\
\midrule
DF4           & $10^{11}$        & 80  & $0.018$ \\
FCC 224       & $3\times10^{10}$ & 70  & $0.007$ \\
NGC 5846-UDG1 & $5\times10^{11}$ & 100 & $0.058$ \\
Dragonfly 44  & $10^{13}$        & 500 & $0.046$ \\
\bottomrule
\end{tabular}
\end{center}

$g_{\rm ext}/a_0\ll 1$ for all four: EFE screening is minor and does
not numerically discriminate DF4 from the matched DM-rich control
NGC~5846-UDG1 (which in fact has a somewhat larger $g_{\rm ext}/a_0$,
not smaller). EFE is therefore rejected as the explanatory mechanism
for DF4's Newton-like kinematics. This is a benchmark-level
observation; the current ECT closure [A4] does not contain an explicit
$g_{\rm ext}$ dependence, so the MOND-EFE comparison is strictly
outside the current ECT formulation.

\section{Honest classification of status}

\begin{description}
\item[(I) Supported by the current preprint.]
\begin{itemize}
\item Closure [A4] has exactly two asymptotic regimes (Newton,
deep-MOND); no third low-$g$ Newtonian branch.
\item Applied to normal RAR galaxies the method returns
$\Xi_{\rm req}\sim O(1)$.
\item Applied to DF4 and FCC~224 the method returns
$\Xi_{\rm req}\ll 1$ robustly under nuisance variation.
\item The Level-B density-only $\Xi$-ansatz [A6] is structurally
insufficient for the DF4 vs NGC~5846-UDG1 matched pair.
\end{itemize}

\item[(II) Reduced-phenomenology upgrades (not yet carried out fully).]
\begin{itemize}
\item Full Jeans with S\'ersic deprojection and aperture matching,
replacing the Wolf proxy.
\item Posterior $P(\Xi\mid\text{data})$ replacing central-value numbers.
\item Same-tracer, same-aperture discipline for matched-control logic.
\end{itemize}

\item[(III) Candidate theory-level directions (none derived in preprint).]
\begin{itemize}
\item Dynamical $G_{\rm eff}(\phi)$ in [A2] could in principle supply
an additional suppression route, but no derivation is given in the
attached files that maps this degree of freedom onto the present
galactic closure for DF4/FCC~224.
\item Non-equilibrium time-dependent $\phi$-dynamics (M1): requires
$c_\ast\ll c$ near criticality, which is not in the current preprint.
\item Orientation-sector coupling via $\alpha_{\rm ECT}$ (M2): would
require a new derivation at galactic scale.
\item History-dependent extensions of $\Xi$ (M3) within the broader
class [\S23722]: requires OP3 resolution.
\item Topological/instanton memory [\S31972]: not connected to the
galactic sector in the preprint.
\end{itemize}
None of (III) is promoted to ``rescue'' status here.

\item[(IV) Cannot be claimed.]
\begin{itemize}
\item DF4/FCC~224 are \emph{not} explained by the current closure.
\item No rescue mechanism is yet derived from first principles in the
attached preprint.
\item The matched-control comparison is a structural challenge, not a
formal impossibility theorem.
\item ``$\bar\phi_{\rm local}\to 0^-\Rightarrow g^\dagger_{\rm eff}\to 0
\Rightarrow$ Newton'' is \emph{not} a derived consequence of the
attached preprint.
\end{itemize}
\end{description}

\section{Future-work factors (secondary, not essential here)}

The following are second-order for the current stress-test diagnosis;
they improve precision of $\Xi_{\rm req}$ or enable more ambitious
analyses, but none alters the qualitative conclusion. Listed as
reminders for future work.

\begin{itemize}
\item Rotation in ``pressure-supported'' dwarfs (DF2 has measurable rotation).
\item GC tracer survival bias and orbital-family selection effects.
\item Mass-sheet degeneracy in external-field analyses.
\item Preprint's 5th-force $L_5=\beta\bar\Psi\gamma^A n_A\Psi$
contribution (Planck-suppressed at galactic scale).
\item Time-since-event correlation with $\Xi_{\rm req}$ via stellar
ages.
\item Spatial $\Xi(R)$ variation within single galaxy (IFU test).
\item Joint ensemble posterior across the whole UDG sample.
\item CMB / early-universe cross-checks for any long-memory mechanism.
\item Interloper mixture modeling in small-$N$ tracer samples.
\item Proper 3D HI cube modeling for AGC~114905.
\item $\Lambda$CDM partial DM halo survival as direct competitor model.
\item Stellar age-metallicity degeneracy in $M/L$.
\item IMF universality assumption.
\end{itemize}

\section{Integration proposal for the preprint (minimal)}

This document is \emph{internal working notes}, not a subsection for
publication. Recommended integration is minimal and is specified in
the companion document {\tt ECT\_UDG\_integration\_proposal.tex}:
\begin{itemize}
\item one row in {\tt tab:future\_discriminants} (entry ND-13);
\item one short paragraph (~150 words) in the galactic-closure
section;
\item \textbf{no} standalone subsection;
\item \textbf{no} claim of derived rescue;
\item \textbf{no} claim that $G_{\rm eff}(\phi)$ solves DF4;
\item \textbf{no} advertisement of the present working notes as a
semi-official supplement to the preprint.
\end{itemize}

\section{Conclusion}

\begin{quote}
The attached preprint already contains a consistent $\phi$-first
galactic closure. DF4 and probably FCC~224 remain sharp stress tests
of that current Level-B closure, not explained consequences of it.
Candidate theory-level routes worth investigating include dynamical
$G_{\rm eff}(\phi)$, non-quasi-static $\phi$-dynamics,
orientation-sector effects, and history-dependent extensions of $\Xi$;
none is presently derived for the galactic closure. The right
integration into the preprint is as an open future discriminant
linked to OP3, not as an explained consequence of the theory.
\end{quote}

\bibliographystyle{unsrt}
\bibliography{references}

\end{document}
